\documentclass[11pt,a4paper]{article}
\usepackage[margin=2cm]{geometry}
\usepackage{booktabs}
\usepackage{array}
\usepackage{longtable}
\usepackage{pdflscape}
\usepackage{hyperref}
\renewcommand{\thetable}{S\arabic{table}}

\begin{document}
\begin{landscape}
\footnotesize
\setlength{\tabcolsep}{3pt}
\setlength{\LTcapwidth}{22cm}
\begin{longtable}{>{\raggedright\arraybackslash}p{3.0cm}>{\raggedright\arraybackslash}p{2.5cm}>{\raggedright\arraybackslash}p{3.7cm}>{\raggedright\arraybackslash}p{4.4cm}>{\raggedright\arraybackslash}p{5.9cm}>{\raggedright\arraybackslash}p{3.3cm}>{\raggedright\arraybackslash}p{1.3cm}}
\caption{The 35 implemented tasks: domain, inspired paradigm, principal
outputs, difficulty mechanism, software verification, and psychometric status.
Scoring fixture (F) = hand-derived numerical fixture for the task's scoring function in the
frontend Vitest suite; Scoring fixture (B1--B4) = hand-derived numerical fixture in backend
pytest batch 1--4; ``+ behavioural tests'' = additional tests of trial generation, response
classification, and difficulty adjustment. A fixture checks a specific calculation, not the
validity of the task. NE = not established (no reliability, validity, or normative data).
Levels are developer-specified; no patients with multiple sclerosis were evaluated.}\label{tab:task-inventory}\\
\toprule Task & Domain & Inspired paradigm & Principal outputs & Difficulty / adaptation & Software verification & Psycho-metric status \\
\midrule
\endfirsthead
\toprule Task & Domain & Inspired paradigm & Principal outputs & Difficulty / adaptation & Software verification & Psycho-metric status \\
\midrule
\endhead
N-Back & Working memory & Visual $n$-back updating~\cite{R1} & accuracy, mean RT, misses, false alarms & Baseline: 20 trials at 1-back; training mode uses 1--3-back by level. Dual N-Back is a separate training task & Scoring fixture (F) & NE \\
Digit Span & Working memory & Digit recall (forward, backward, mixed)~\cite{R2} & trial accuracy, span, score & Levels 1--10: length 3--9 and forward/backward/mixed recall; trials 4 and 7 use one level higher & Scoring fixture (B1) & NE \\
Spatial Span & Working memory & Spatial-sequence recall (Corsi-type)~\cite{R3} & set and order accuracy, maximum span, score & Levels 1--10: sequence length, display and interval timing & Scoring fixture (B2) & NE \\
Letter--Number Sequencing & Working memory & Ordered manipulation of mixed letters and digits~\cite{R2} & set accuracy, order accuracy, speed, score & Levels 1--10: sequence composition and length, timing & Scoring fixture (B2) & NE \\
Operation Span & Working memory & Complex span (arithmetic check plus letter recall)~\cite{R4} & arithmetic accuracy, letter accuracy, dual-task score & Levels 1--10: set size, operation complexity, timing & Scoring fixture (B2) & NE \\
Dual N-Back & Working memory & Dual-modality (position and letter) $n$-back~\cite{R5} & visual, auditory and dual accuracy, score & Levels 1--10: $n=1$--3, stream length, timing, target rate & Scoring fixture (B4) + behavioural tests & NE \\
\midrule
Simple Reaction Time & Processing speed & Simple reaction time~\cite{R6} & mean RT, RT SD, choice accuracy, combined score & Shared processing-speed workflow; levels 1--10 change delay, alternatives, timing & Scoring fixture (F) & NE \\
Symbol--Digit Matching & Processing speed & Symbol--digit substitution~\cite{R7} & correct, incorrect, completion time, score & Levels 1--10: symbol count, item count, exposure, target & Scoring fixture (B1) & NE \\
Numeric Trail Task & Processing speed & Sequential visual trail~\cite{R8} & completion time, errors, score & Levels 1--10: node count, layout, expected time & Scoring fixture (B2) & NE \\
Pattern Comparison & Processing speed & Rapid same/different comparison~\cite{R9} & accuracy, RT, timeouts, score & Levels 1--10: grid size, similarity, changed cells, duration, trials & Scoring fixture (B1) & NE \\
Inspection Time & Processing speed & Brief line discrimination with mask~\cite{R10} & accuracy, threshold duration, score & Levels 1--10: presentation time, line difference, trials & Scoring fixture (B1) & NE \\
Choice Reaction Time & Processing speed & Stimulus--response choice RT~\cite{R11} & accuracy, mean RT, consistency, score & Levels 1--10: 2--4 alternatives, trials, windows; up if score $\ge83$ and accuracy $\ge82$, down if $<63$ & Scoring fixture (B4) + behavioural tests & NE \\
\midrule
Continuous Performance Task & Attention & Sustained target detection (AX-type baseline)~\cite{R12} & target hit rate, false clicks, mean RT, score & Baseline: 60 trials; training mode adds 10 trials per level above 3. Go/No-Go is a separate training task & Scoring fixture (F) & NE \\
Serial Addition & Attention & Serial addition (PASAT-type)~\cite{R13} & correct, omissions, RT, accuracy-based score & Levels 1--10: shorter inter-stimulus interval, trials & Scoring fixture (B1) & NE \\
Color--Word Interference & Attention & Color--word interference~\cite{R14} & accuracy, interference score, RT, score & Levels 1--10: color set, congruency ratio, timing, trials & Scoring fixture (B3) & NE \\
Go/No-Go & Attention & Response inhibition~\cite{R30,R31} & go and no-go accuracy, RT, consistency, score & Levels 1--10: stimulus similarity, no-go frequency, timing, trials & Scoring fixture (B3) & NE \\
Flanker & Attention & Flanker interference~\cite{R15} & accuracy, congruent/incongruent RT, conflict control, score & Levels 1--10: congruency ratio, spacing, timing, trials & Scoring fixture (B3) & NE \\
Sustained-Attention Response & Attention & Frequent-go, rare-no-go vigilance~\cite{R16} & overall and no-go accuracy, RT, vigilance penalty, score & Levels 1--10: stream length, target rate, timing; up if score $\ge84$ and $\le4$ commissions, down if $<64$ & Scoring fixture (B4) + behavioural tests & NE \\
\midrule
Task Switching & Flexibility & Cued task switching~\cite{R17} & accuracy, switch cost, perseveration errors, score & Baseline: 40 trials; training mode adds 10 trials per level above 4. Rule Shift is a separate training task & Scoring fixture (F) & NE \\
Alternating Trail Task & Flexibility & Alternating number--letter trail~\cite{R8} & completion time, errors, B-minus-A & Levels 1--10: node count, layout, expected time & Scoring fixture (B3) & NE \\
Card-Sorting Rule Shift & Flexibility & Feedback-guided card sorting~\cite{R18} & categories, perseverative errors, accuracy, score & Levels 1--10: trials, rule-change frequency, feedback, category complexity & Scoring fixture (B3) & NE \\
Dimensional Card Sort & Flexibility & Dimensional rule switching~\cite{R19} & accuracy, mean RT, switch cost, score & Levels 1--10: dimensions, trials, switch frequency, timing & Scoring fixture (B3) & NE \\
Rule Shift & Flexibility & Cued set shifting (color, shape, count)~\cite{R17} & overall and switch accuracy, RT, perseverative errors, score & Levels 1--10: block sequence, cue time; up if score $\ge83$, switch accuracy $\ge72$ and $\le5$ perseverative errors, down if $<63$ & Scoring fixture (B4) + behavioural tests & NE \\
Plus--Minus Switching & Flexibility & Alternating arithmetic switching~\cite{R20} & accuracy, mean RT, switch cost, score & Levels 1--10: operand range, mixed-block frequency, items, timing & Scoring fixture (B3) & NE \\
\midrule
Tower Planning & Planning & Tower-type planning~\cite{R21} & problems solved, moves, planning efficiency, time, score & Levels 1--10: problem count, disks, minimum moves, time limits & Scoring fixture (B3) & NE \\
Stockings-Style Planning & Planning & Tower-style planning with balls~\cite{R22} & problems solved, planning efficiency, mean time, score & Levels 1--10: problems, balls, minimum moves, time limits & Scoring fixture (B3) & NE \\
Letter Fluency & Planning & Phonemic fluency~\cite{R23} & valid unique words, errors, mean words, score & Levels 1--10: letter sets, prompt count, duration & Scoring fixture (B3) & NE \\
Category Fluency & Planning & Semantic fluency~\cite{R23} & unique valid words, duplicates, score & Levels 1--10: category pools; expected count rises with level & Scoring fixture (B4) + behavioural tests & NE \\
Twenty Questions & Planning & Constraint-seeking problem solving~\cite{R24} & identification, questions used, strategy score, score & Levels 1--10: five object pools in rotation; score multiplier $1+0.05(\mathrm{level}-1)$ & Scoring fixture (B4) + behavioural tests & NE \\
\midrule
Visual Search & Visual scanning & Target search among distractors~\cite{R25} & targets found, accuracy, time per target, score & Levels 1--10: item count, distractor similarity, targets, grid & Scoring fixture (B4) + behavioural tests & NE \\
Cancellation & Visual scanning & Letter or symbol cancellation~\cite{R26} & hits, omissions, commissions, time, score & Levels 1--10: 150--540 cells, target count and types & Scoring fixture (B4) + behavioural tests & NE \\
Feature and Conjunction Search & Visual scanning & Feature-integration search~\cite{R25} & accuracy, search time, set-size slope & Levels 1--10: set size, similarity, conjunction load & Scoring fixture (B3) & NE \\
Landmark Line Judgment & Visual scanning & Line-bisection judgment~\cite{R27} & accuracy, omissions, spatial bias, RT, score & Levels 1--10: offset distribution, exposure; up if score $\ge84$, $|\mathrm{bias}|\le8$ and $\le1$ omission, down if $<64$ & Scoring fixture (B4) + behavioural tests & NE \\
Multiple-Object Tracking & Visual scanning & Dynamic object tracking~\cite{R28} & accuracy, precision, recall, F1, score & Levels 1--10: 6--16 objects, 1--3 targets, speed, duration & Scoring fixture (B4) + behavioural tests & NE \\
Central--Peripheral Attention & Visual scanning & Divided central/peripheral attention~\cite{R29} & central accuracy, peripheral localization, RT, score & Levels 1--10 shorten presentation time (800 to 150 ms): levels 1--3 central only, 4--6 add a peripheral target, 7--10 add 24--47 distractors & Scoring fixture (B2) & NE \\
\bottomrule
\end{longtable}

\vspace{0.6em}
\begin{minipage}{0.95\linewidth}\footnotesize
\textbf{Note.} N-Back, Continuous Performance Task, Task Switching, and Simple Reaction Time are scored client-side in the baseline workflow. Their scoring functions are extracted to a frontend module and each is covered by one Vitest fixture (F); the other 31 tasks are scored by backend services and covered by pytest fixtures (B1--B4).
\end{minipage}

\end{landscape}

\clearpage
\footnotesize
\begin{thebibliography}{31}
\bibitem{R1} W.K. Kirchner, Age differences in short-term retention of rapidly changing information, J. Exp. Psychol. 55 (1958) 352--358.
\bibitem{R2} D. Wechsler, Wechsler Adult Intelligence Scale, fourth ed., Pearson, 2008.
\bibitem{R3} R.P.C. Kessels et al., The Corsi Block-Tapping Task: standardization and normative data, Appl. Neuropsychol. 7 (2000) 252--258.
\bibitem{R4} A.R.A. Conway et al., Working memory span tasks: a methodological review and user's guide, Psychon. Bull. Rev. 12 (2005) 769--786.
\bibitem{R5} S.M. Jaeggi et al., Improving fluid intelligence with training on working memory, Proc. Natl. Acad. Sci. USA 105 (2008) 6829--6833.
\bibitem{R6} F.C. Donders, On the speed of mental processes, Acta Psychol. 30 (1969) 412--431 (original work published 1868).
\bibitem{R7} A. Smith, Symbol Digit Modalities Test manual, Western Psychological Services, 1982.
\bibitem{R8} R.M. Reitan, Validity of the Trail Making Test as an indicator of organic brain damage, Percept. Mot. Skills 8 (1958) 271--276.
\bibitem{R9} T.A. Salthouse, R.L. Babcock, Decomposing adult age differences in working memory, Dev. Psychol. 27 (1991) 763--776.
\bibitem{R10} D. Vickers, M.D. Nettelbeck, J. Willson, Perceptual indices of performance: the measurement of inspection time and noise in the visual system, Perception 1 (1972) 263--295.
\bibitem{R11} W.E. Hick, On the rate of gain of information, Q. J. Exp. Psychol. 4 (1952) 11--26.
\bibitem{R12} H.E. Rosvold et al., A continuous performance test of brain damage, J. Consult. Psychol. 20 (1956) 343--350.
\bibitem{R13} D.M.A. Gronwall, Paced auditory serial-addition task: a measure of recovery from concussion, Percept. Mot. Skills 44 (1977) 367--373.
\bibitem{R14} J.R. Stroop, Studies of interference in serial verbal reactions, J. Exp. Psychol. 18 (1935) 643--662.
\bibitem{R15} B.A. Eriksen, C.W. Eriksen, Effects of noise letters upon identification of a target letter, Percept. Psychophys. 16 (1974) 143--149.
\bibitem{R16} I.H. Robertson et al., ``Oops!'': performance correlates of everyday attentional failures, Neuropsychologia 35 (1997) 747--758.
\bibitem{R17} R.D. Rogers, S. Monsell, Costs of a predictable switch between simple cognitive tasks, J. Exp. Psychol. Gen. 124 (1995) 207--231.
\bibitem{R18} E.A. Berg, A simple objective technique for measuring flexibility in thinking, J. Gen. Psychol. 39 (1948) 15--22.
\bibitem{R19} P.D. Zelazo, The Dimensional Change Card Sort, Nat. Protoc. 1 (2006) 297--301.
\bibitem{R20} A.T. Jersild, Mental set and shift, Arch. Psychol. 14 (1927) 1--81.
\bibitem{R21} T. Shallice, Specific impairments of planning, Philos. Trans. R. Soc. Lond. B 298 (1982) 199--209.
\bibitem{R22} A.M. Owen et al., Planning and spatial working memory following frontal lobe lesions, Neuropsychologia 28 (1990) 1021--1034.
\bibitem{R23} A.L. Benton, Differential behavioral effects in frontal lobe disease, Neuropsychologia 6 (1968) 53--60.
\bibitem{R24} F.A. Mosher, J.R. Hornsby, On asking questions, in: J.S. Bruner et al. (Eds.), Studies in Cognitive Growth, Wiley, 1966.
\bibitem{R25} A.M. Treisman, G. Gelade, A feature-integration theory of attention, Cogn. Psychol. 12 (1980) 97--136.
\bibitem{R26} S. Weintraub, M.-M. Mesulam, Mental state assessment of young and elderly adults, in: Principles of Behavioral Neurology, F.A. Davis, 1985.
\bibitem{R27} G. Jewell, M.E. McCourt, Pseudoneglect: a review and meta-analysis, Neuropsychologia 38 (2000) 93--110.
\bibitem{R28} Z.W. Pylyshyn, R.W. Storm, Tracking multiple independent targets, Spat. Vis. 3 (1988) 179--197.
\bibitem{R29} K. Ball et al., Visual attention problems as a predictor of vehicle crashes in older drivers, Invest. Ophthalmol. Vis. Sci. 34 (1993) 3110--3123.
\bibitem{R30} D.J. Simmonds, J.J. Pekar, S.H. Mostofsky, Meta-analysis of Go/No-go tasks demonstrating that fMRI activation associated with response inhibition is task-dependent, Neuropsychologia 46 (2008) 224--232.
\bibitem{R31} A. Diamond, Executive functions, Annu. Rev. Psychol. 64 (2013) 135--168.
\end{thebibliography}
\normalsize

\end{document}
